\documentclass[12pt,a4paper]{article}

\usepackage{geometry}
\geometry{margin=1.2in}

\usepackage{amsmath}
\usepackage{booktabs}
\usepackage{hyperref}
\usepackage{microtype}
\usepackage{graphicx}
\usepackage{longtable}
\usepackage{array}
\usepackage{xcolor}

\title{Preliminary Analysis of Developer Activity and Language Ecosystems\\
in a Large GitHub Follow Network}
\author{Flyxion}
\date{\today}

\begin{document}

\maketitle

\begin{abstract}
This report presents preliminary findings from a large-scale observational study
of developer activity within a GitHub follow network. Using the GitHub GraphQL
API, we collect activity metadata for accounts followed by a single user whose
follow graph exceeds one million accounts. Each account is classified as active
or inactive based on the recency of pushes to public repositories, and the
dominant programming language of the most recently updated repository is
recorded.

Analysis of the first $\sim150{,}000$ accounts reveals systematic differences in
activity levels across programming language ecosystems. The overall active
fraction is approximately $0.589$. Newer ecosystems such as Astro, Rust, and
TypeScript exhibit activity ratios above $0.83$, while older ecosystems such as
Ruby and Java fall near $0.44$--$0.51$. Unexpectedly, the developer survival
curve is not monotonically decreasing: active fractions decline from account ages
zero through eleven years, reaching a minimum near $0.52$, and then \emph{rise}
steadily for accounts aged twelve to seventeen years, recovering to fractions
above $0.66$. This U-shaped pattern suggests that the oldest GitHub accounts in
the dataset are disproportionately held by developers who have remained
continuously active since the platform's earliest years. These results provide an
empirical snapshot of the structure and generational stratification of open
source developer communities.
\end{abstract}

\newpage
\section{Introduction}

GitHub has become one of the central platforms for collaborative software
development, and its public data provides a unique opportunity to study the
dynamics of programming language ecosystems and developer activity over time.

Most existing studies rely on repository-level datasets such as GitHub Archive
or BigQuery mirrors of repository metadata. While these datasets capture
repository events in detail, they provide only indirect insight into the life
cycles of individual developers.

The present study instead analyzes a \emph{developer-centered dataset}. Starting
from a single GitHub account that follows more than one million other accounts,
we crawl the follow network using the GitHub GraphQL API and collect basic
metadata describing each followed account. The goal is to estimate:
\begin{itemize}
\item the fraction of developers who remain active over time,
\item differences in activity rates across programming language ecosystems,
\item and the approximate lifetime of developer participation on GitHub.
\end{itemize}

The results reported here are preliminary and based on the first
$\sim150{,}000$ accounts processed during an ongoing crawl.

\section{Dataset Construction}

The dataset is collected using the GitHub GraphQL API. For each followed account
we retrieve:
\begin{itemize}
\item account creation timestamp,
\item number of followers,
\item number of public repositories,
\item the most recently pushed public repository,
\item and the primary language of that repository.
\end{itemize}

Accounts are classified as \textbf{active} if the most recent repository push
occurred within the previous year; otherwise the account is classified as
\textbf{inactive}. Accounts for which no language-tagged repository is available
appear in the dataset under the label \textsc{None}.

For the $N = 154{,}283$ accounts processed at the time of this report, the
dataset contains:

\begin{center}
\begin{tabular}{lr}
\toprule
Active accounts   & 90,808 \\
Inactive accounts & 63,475 \\
\midrule
Active fraction   & $\approx 0.589$ \\
\bottomrule
\end{tabular}
\end{center}

Because the sample is drawn from a follow network rather than uniformly across
GitHub, it is biased toward more visible and historically active developers.
This selection effect is discussed further in Section~\ref{sec:survival}.

\section{Language Activity Ratios}

For each programming language we compute an \emph{activity ratio}
\[
R \;=\; \frac{\text{active}}{\text{active}+\text{inactive}}.
\]
This metric estimates the probability that a developer whose most recently
active repository is tagged with a given language is still committing code.
Table~\ref{tab:activity} reports results for all languages with at least 336
accounts in the dataset, ordered by activity ratio.

\begin{longtable}{lrrrr}
\toprule
Language & Active & Inactive & Total & $R$ \\
\midrule
\endhead
\midrule
\multicolumn{5}{r}{\small\itshape continued on next page}\\
\endfoot
\bottomrule
\endlastfoot
Astro              &    306 &     39 &     345 & 0.89 \\
MDX                &    303 &     48 &     351 & 0.86 \\
Rust               &  2160 &    400 &    2560 & 0.84 \\
TypeScript         & 11988 &   2413 &   14401 & 0.83 \\
Lua                &    692 &    169 &     861 & 0.80 \\
Dart               &    588 &    221 &     809 & 0.73 \\
Go                 &  2460 &    991 &    3451 & 0.71 \\
Python             & 16890 &   6973 &   23863 & 0.71 \\
Shell              &  2242 &   1011 &    3253 & 0.69 \\
Kotlin             &    696 &    326 &    1022 & 0.68 \\
HCL                &    239 &    119 &     358 & 0.67 \\
HTML               &  6617 &   3320 &    9937 & 0.67 \\
SCSS               &    401 &    200 &     601 & 0.67 \\
PowerShell         &    237 &    126 &     363 & 0.65 \\
Dockerfile         &    217 &    119 &     336 & 0.65 \\
TeX                &    284 &    166 &     450 & 0.63 \\
Swift              &    565 &    331 &     896 & 0.63 \\
R                  &    677 &    421 &    1098 & 0.62 \\
Vue                &    514 &    329 &     843 & 0.61 \\
JavaScript         &  8531 &   5357 &   13888 & 0.61 \\
Jupyter Notebook   &  3749 &   2486 &    6235 & 0.60 \\
C++                &  2757 &   1822 &    4579 & 0.60 \\
C\#                &  1444 &   1075 &    2519 & 0.57 \\
C                  &  1619 &   1222 &    2841 & 0.57 \\
CSS                &  1227 &   1008 &    2235 & 0.55 \\
PHP                &  1154 &    987 &    2141 & 0.54 \\
Java               &  2649 &   2502 &    5151 & 0.51 \\
MATLAB             &    189 &    187 &     376 & 0.50 \\
Ruby               &    674 &    848 &    1522 & 0.44 \\
\textsc{None}      & 14682 &  26193 &   40875 & 0.36 \\
\caption{Activity ratios for languages with at least 336 accounts ($N=154{,}283$).
Ordered by activity ratio descending; \textsc{None} placed last.}
\label{tab:activity}
\end{longtable}

Several patterns are immediately visible.

Newer language ecosystems lead the table. Astro ($R = 0.89$) and MDX ($R =
0.86$) represent communities whose adoption is almost entirely post-2020;
virtually every developer who has ever used these tools remains active by
construction. Rust ($0.84$) and TypeScript ($0.83$) reflect the strong
contemporary growth of systems-level and typed JavaScript development
respectively.

At the lower end, Ruby ($0.44$) and Java ($0.51$) carry large populations of
historical accounts from their respective peaks of GitHub adoption. MATLAB
($0.50$) sits near the overall median, consistent with a community of academic
and engineering users whose engagement with public version control is more
episodic than that of professional software developers.

The \textsc{None} category ($R = 0.36$) is substantially below the overall mean
($0.589$). This gap reflects two overlapping effects: inactive developers are
less likely to maintain language-tagged repositories, and accounts with no public
repositories at all---including spam, bot, and promotional accounts---are
counted here. Further decomposition of this category is deferred to future work.

\subsection{Compositional Asymmetry Between Active and Inactive Populations}

The active and inactive populations differ not only in aggregate size but in
their language composition. Among active accounts, Python (18.6\%), TypeScript
(13.2\%), and JavaScript (9.4\%) are the three most represented languages. Among
inactive accounts, \textsc{None} accounts for 41.3\% of the population---more
than twice its share among active accounts (16.2\%). TypeScript, by contrast,
constitutes 13.2\% of active accounts but only 3.8\% of inactive accounts,
illustrating the asymmetric concentration of active development in newer
ecosystems.

\subsection{The Objective-C Signal}

Objective-C does not appear in Table~\ref{tab:activity} at the 150K snapshot
because its sample size has not yet crossed the reporting threshold in this
iteration of the analysis. At the 30K preliminary snapshot it exhibited $R =
0.15$ (9 active out of 59 total), the lowest ratio in the dataset by a
substantial margin. This figure is consistent with the wholesale migration of
Apple platform development from Objective-C to Swift following Swift's
introduction at WWDC 2014 and its open-sourcing in 2015. Developers whose most
recently active repository is still tagged Objective-C constitute a population
skewed heavily toward inactivity: those who have not migrated to Swift are, by
and large, those who have stopped contributing publicly altogether. The ratio
functions as a fossil record of a completed platform transition rather than a
measure of an ongoing ecosystem.

\section{Account Age and Developer Survival}\label{sec:survival}

Using account creation timestamps and most recent push timestamps, we estimate
the fraction of developers remaining active as a function of account age.
Table~\ref{tab:survival} presents the full survival curve across eighteen annual
bins.

\begin{table}[h]
\centering
\begin{tabular}{rc}
\toprule
Account Age (years) & Active Fraction \\
\midrule
 0 & 0.631 \\
 1 & 0.656 \\
 2 & 0.648 \\
 3 & 0.620 \\
 4 & 0.608 \\
 5 & 0.609 \\
 6 & 0.597 \\
 7 & 0.574 \\
 8 & 0.556 \\
 9 & 0.540 \\
10 & 0.527 \\
11 & 0.524 \\
12 & 0.538 \\
13 & 0.565 \\
14 & 0.606 \\
15 & 0.633 \\
16 & 0.661 \\
17 & 0.680 \\
18 & 0.632 \\
\bottomrule
\end{tabular}
\caption{Active fraction by account age bin ($N=154{,}283$).}
\label{tab:survival}
\end{table}

\subsection{A U-Shaped Survival Curve}

The survival curve in Table~\ref{tab:survival} is not monotonically decreasing,
as a simple dropout model would predict. Instead it exhibits a pronounced
U-shaped form: the active fraction declines from approximately $0.631$ at age
zero through a minimum of $0.524$ at age eleven, then \emph{rises steadily} from
age twelve onward, reaching $0.680$ at age seventeen before a partial decline at
age eighteen.

This pattern is unexpected and warrants careful interpretation. We propose the
following account.

GitHub was founded in April 2008. Accounts that are now aged twelve to seventeen
years were therefore created between approximately 2009 and 2014---the platform's
earliest adoption cohort. This population consists disproportionately of
developers who were already deeply embedded in open source culture at the time of
GitHub's founding: early adopters, maintainers of established projects, and
professionals for whom public version control was a durable practice rather than
a transient experiment. Such developers are precisely those most likely to have
sustained activity across a fifteen-year period. In other words, the oldest
accounts that are still present in a follow network are not a random sample of
early GitHub users: they are the \emph{survivors}---a highly selected
subpopulation of long-term contributors who were never likely to go inactive.

The intermediate cohort (ages five to eleven, corresponding to account creation
roughly 2015--2021) catches the broad wave of GitHub adoption that followed the
platform's mainstream growth. This population is more heterogeneous: it includes
students, hobbyists, and developers who created accounts for course projects or
brief professional forays, many of whom have since become inactive. The trough in
the survival curve at ages ten to eleven thus reflects the high churn
characteristic of this mass-adoption wave.

The youngest cohort (ages zero to two, accounts created 2024--2026) shows a
moderately high active fraction ($0.63$--$0.66$), but this figure is partially
an artifact of left-censoring: accounts created recently and already gone
inactive are less likely to appear in a follow network, because inactive accounts
accumulate fewer followers over time.

Taken together, the U-shaped curve is better understood as a \emph{selection
artifact} than as a true survival phenomenon. It reflects the varying composition
of different account-age cohorts as filtered through the follow-network sampling
frame, rather than a literal rise in individual developer activity with age.

\subsection{Absence of a Simple Half-Life}

The preliminary 30K analysis suggested an approximate developer activity
half-life of ten years, defined as the account age at which the active fraction
first crosses $0.5$. The 150K data substantially revise this picture: the active
fraction does not cross $0.5$ at any age bin in the current dataset. The minimum
observed value is $0.524$ at age eleven. This reflects the follow-network
selection bias noted above: the seed account's follow list is not a
representative cross-section of all GitHub accounts, but a set of accounts
that attracted sufficient attention to be followed, which selects for persistent
activity. A simple half-life estimate derived from this curve would not be
interpretable as a property of the general developer population.

\section{Language-Specific Developer Lifetimes}

We compute median observed developer lifetimes by language, defined as the
median interval between account creation and most recent repository push over all
accounts in the language with both timestamps present.
Table~\ref{tab:lifetime} presents results for the twenty languages with the
longest median lifetimes among those with at least 20 observations.

\begin{table}[h]
\centering
\begin{tabular}{lrr}
\toprule
Language & Median Lifetime (years) & $n$ \\
\midrule
Erlang       & 12.65 &   22 \\
Starlark     & 12.54 &   25 \\
Common Lisp  & 11.75 &   53 \\
Clojure      & 11.75 &  169 \\
Elixir       & 11.62 &  163 \\
Emacs Lisp   & 11.52 &  184 \\
Go Template  & 10.90 &   53 \\
OCaml        & 10.82 &   46 \\
Haskell      & 10.63 &   98 \\
Vim Script   & 10.55 &  257 \\
Zig          & 10.50 &   81 \\
Scheme       & 10.48 &   20 \\
LLVM         & 10.37 &   43 \\
Jinja        &  9.67 &   66 \\
Nix          &  9.65 &  258 \\
Nunjucks     &  9.55 &   36 \\
Perl         &  9.46 &  162 \\
F\#          &  9.40 &   28 \\
Rust         &  9.18 & 2560 \\
Just         &  9.18 &   24 \\
\bottomrule
\end{tabular}
\caption{Top 20 languages by median developer lifetime ($n \geq 20$).}
\label{tab:lifetime}
\end{table}

The languages at the top of Table~\ref{tab:lifetime} form a coherent profile:
functional and systems languages (Erlang, Common Lisp, Clojure, Elixir, OCaml,
Haskell, Scheme, F\#), editor-specific scripting languages (Emacs Lisp, Vim
Script), and languages associated with infrastructure and build tooling
(Starlark, Go Template, Nix, LLVM, Just). These communities attract developers
with deep technical commitments who maintain long-term engagement with their
tools regardless of broader industry trends.

Zig ($10.50$ years, $n=81$) is a notable exception: a young language (publicly
released 2016) whose users nonetheless exhibit long measured lifetimes. This
reflects the fact that Zig attracted experienced systems programmers who already
held old GitHub accounts, rather than new entrants to the platform.

The contrast between median lifetime and activity ratio remains instructive at
the 150K snapshot. Rust has both a long median lifetime ($9.18$ years, $n =
2{,}560$) and a high activity ratio ($0.84$), making it unusual: a large
community that is both historically grounded and highly retained. Clojure and
Elixir exhibit long lifetimes ($11.75$ and $11.62$ years respectively) but more
moderate activity ratios, consistent with stable, slowly growing communities of
committed practitioners.

\section{Discussion}

The 150K snapshot substantially extends and in some cases revises the picture
drawn from the 30K preliminary data.

The language activity ratios have stabilized: the rank ordering of major
languages is consistent with the earlier snapshot, and the sample sizes for the
largest languages (Python at 23,863; TypeScript at 14,401; JavaScript at 13,888)
are sufficient to treat these ratios as reliable estimates of population-level
activity rates. Small languages with fewer than a few hundred observations remain
subject to significant sampling variability and should be interpreted with
caution.

The survival curve is the most substantively novel finding at this scale. The
U-shaped form---and in particular the recovery of active fractions among accounts
aged twelve to seventeen years---cannot be explained by a simple model of
developer dropout. It requires acknowledging the compositional heterogeneity of
the population: the oldest surviving accounts in a follow network are not typical
early users, but a selected cohort of long-term contributors. This selection
effect is intrinsic to the sampling frame and will persist in the full dataset.
Correcting for it would require an estimate of the follower-count distribution
conditional on account age and activity status, which the full dataset will
make possible.

\section{Future Work}

The present analysis covers approximately 15\% of the full follow list (50,040
batches total, 1,000,800 accounts). Completion of the crawl will enable:
\begin{itemize}
\item precise activity ratio estimates for smaller language communities currently
      below the reporting threshold,
\item formal survival modeling (Weibull, log-normal, or mixture models) to
      characterize attrition across cohorts,
\item disaggregated survival curves by language and account vintage to test
      whether the U-shaped pattern is uniform or language-specific,
\item and analysis of the relationship between follower count and long-term
      activity, which would quantify the follow-network selection bias directly.
\end{itemize}

Because the dataset is developer-centered rather than repository-centered, it
provides a perspective on open source community dynamics not available from
existing large-scale GitHub datasets.

\section{Crawl Architecture and Analysis Pipeline}

\subsection{Data Collection}

The crawler issues paginated requests to the GitHub GraphQL API, traversing the
follower list of a seed account whose follow graph exceeds one million accounts.
Requests are batched at exactly 20 accounts per query. The total follow list
spans approximately 1,000,800 accounts distributed across 50,040 batches; at the
time of this report, approximately 7,700 batches ($\sim15\%$) had been completed.
Results are written incrementally to three output files:
\texttt{following-activity.csv} (the primary structured record),
\texttt{active-following.txt}, and \texttt{inactive-following.txt}, permitting
intermediate analyses to be run against any partial snapshot without interrupting
collection.

The CSV schema exposes the following fields per account:

\begin{center}
\begin{tabular}{ll}
\toprule
Field & Description \\
\midrule
\texttt{status}            & \texttt{ACTIVE} or \texttt{INACTIVE} \\
\texttt{primary\_language} & Language tag of most recently pushed repository \\
\texttt{created\_at}       & Account creation timestamp (ISO\,8601) \\
\texttt{last\_pushed\_at}  & Timestamp of most recent public repository push \\
\bottomrule
\end{tabular}
\end{center}

An account is labeled \texttt{ACTIVE} if \texttt{last\_pushed\_at} falls within
the twelve months preceding the crawl date; otherwise it is labeled
\texttt{INACTIVE}. The crawler implements retry logic with attempt tracking,
allowing transient API errors to be handled without restarting the full
traversal.

\subsection{Analysis Scripts}

Three Python analysis scripts operate on the CSV output, invoked together by a
shell driver (\texttt{run\_analysis.sh}) that tees each script's output to a
corresponding file in an \texttt{analysis-results/} directory.

\paragraph{\texttt{analyse\_activity.py}} Computes the overall active/inactive
split, per-language activity ratios (filtered to languages with at least 100
total accounts, sorted by total presence), and a year-resolution survival table
derived from \texttt{created\_at} timestamps. The approximate half-life is read
from the survival table as the first bin whose active fraction falls below $0.5$;
as discussed in Section~\ref{sec:survival}, this threshold is not crossed in the
current dataset.

\paragraph{\texttt{analyse\_languages.py}} Separately tabulates the top languages
by count among active and inactive accounts, reporting each language's share of
its respective subpopulation. This decomposition exposes compositional
asymmetries not visible in the combined activity ratio.

\paragraph{\texttt{language\_survival.py}} Computes per-language median developer
lifetimes as the median of
$\{(\texttt{last\_pushed\_at} - \texttt{created\_at}) / 365.25\}$
over all accounts with both timestamps present, excluding languages with fewer
than 20 observations.

\subsection{Rate Limiting and Projected Completion}

At 20 accounts per GraphQL batch, the full follow list of $\sim1{,}000{,}800$
accounts requires 50,040 batches. GitHub imposes a rate limit of 5,000 GraphQL
points per hour on authenticated requests; the effective crawl rate is governed
by the point cost of each batch query. The incremental write design ensures that
each intermediate snapshot is self-consistent: the active/inactive classification
threshold is applied at analysis time against the stored
\texttt{last\_pushed\_at} field, so the definition of activity remains stable
across repeated analyses of different snapshots.

\section{Conclusion}

Analysis of 154,283 accounts from a large GitHub follow network reveals
systematic variation in developer activity across programming language ecosystems
and across account-age cohorts. Astro, MDX, Rust, and TypeScript lead with
activity ratios of $0.83$--$0.89$; Ruby, Java, and the no-language
(\textsc{None}) category trail at $0.36$--$0.51$. The developer survival curve
is U-shaped rather than monotonically decreasing, with the oldest accounts (aged
12--17 years) showing higher active fractions than middle-aged accounts, a
pattern attributable to survivor selection within the follow-network sampling
frame. Median developer lifetimes are longest in functional and infrastructure
language communities and shorter in mass-adoption ecosystems.

Completion of the full one-million-account dataset will enable formal survival
modeling, disaggregated cohort analysis, and direct quantification of the
follow-network selection bias.

\end{document}
